% BEGIN PROFESSOR FEEDBACK PF25: Discussion - Detection value
\subsection{Detection value of temporal motifs}

RQ1 concerns the complete fingerprint representation. The comparison tests
whether adding local motif-role features to the 17 conventional features improves
detection. Since the roles combine structure, direction, focal-node position,
and edge order, the result does not isolate a causal contribution of temporal
ordering.

The first result is that temporal motifs contain a useful signal, but that their
additional value depends on which other features are already present. Adding the
motif fingerprint to the baseline increases AP in all six HuggingFace
comparisons. Three of these increases are significant. The non-significant
comparisons include the prevalence cohort of one 32-role sample and both
cohorts of the 40-role analysis. In that analysis, the improvements are 0.0237 and
0.0161, but both confidence intervals include zero. The
binary Real-CATS comparison is even closer, with an AP difference of only
0.0011. It also produces a lower F1 score after the local motif-role features are added.
The positive direction is therefore not enough to conclude that motifs provide
a reliable improvement over the baseline.

This does not mean that the motif fingerprint contains no information. The
motif-only models can still rank addresses to some degree, and the Real-CATS
motif-only model reaches an AP of 0.8708. More importantly, the
role-by-role fingerprint performs better than total motif volume alone. In the
40-role HuggingFace analysis, this difference is 0.1173 AP in the controlled
cohort and 0.1528 AP in the prevalence cohort. This shows that the distribution
of activity across motif roles is more informative than the number of motifs by
itself.

The result changes once conventional features are added. A baseline supplemented
with motif composition and motif volume performs similarly to the full combined
model across all three HuggingFace samples. The detailed role counts therefore
contain a signal,
but much of that signal overlaps with information that can already be obtained
from simpler summaries. This overlap is plausible because many transaction
patterns affect several feature groups at the same time. A rapid fan-out, for
example, changes the motif fingerprint, but it can also increase the number of
outgoing transfers, the out-degree, and the number of transaction partners. The
comparison does not isolate the exact source of this overlap, but it shows that
the complete fingerprint adds little once these related properties are present.

The triangle analysis leads to the same interpretation. Extending the
HuggingFace graph makes it possible to retain all eight triangle roles, but
adding them changes AP by at most 0.0071 in the reported comparisons. None of
these differences is significant. Triangle motifs may still be useful in other
graphs or for other kinds of malicious behaviour. The present results only show
that they do not strengthen address detection in these cohorts.

The strong baseline performance is consistent with earlier Ethereum detection
work based on transaction aggregates and static graph features
\citep{chen2022phishing, sanjalawe2023atd}. The motif-only results also agree
with the observation of \citet{arnold2024motifs} that local motif signatures can
expose behavioural differences that disappear in global counts. This thesis
tests the next step by asking whether that information improves a detector on
held-out addresses. The comparison shows that a descriptive signal does not
automatically provide additional predictive value. The result
also differs from the improvement reported by \citet{liu2023temporal} on
non-blockchain financial networks. This suggests that results from one type of
transaction network do not necessarily transfer to another, especially when
the labels and baseline features are different.

Temporal motif features have also been used for Ethereum phishing detection
\citep{wang2023temporal}. The present results qualify the usefulness of this
representation under the labels and evaluation protocol used here. They do not
provide a direct comparison with that earlier detector.
% END PROFESSOR FEEDBACK PF25
